\begin{lemma}[Cyclic $\Omega_u$-contraction for blocked sector gauges]%
    \label{lem:sector_tensor_proportional_blocked_match}
    \lean{MPSTensor.sectorTensor_proportional_of_blockedMatch}
    \leanok
    \uses{lem:sector_match_propagation, def:cyclic_sector_decomp, def:periodic_corner_letter, def:periodic_corner_product, lem:periodic_corner_product_append, lem:periodic_corner_product_blockmatch_pow, lem:periodic_corner_contraction, thm:periodic_product_one_phase_extraction, thm:left_canonical_sector_scalar_normalization, lem:finite_cycle_coboundary}
    Let $A$ and $B$ be left-canonical tensors with period $m$, and let
    $\bar A\sim\bigoplus_u A_u$ and $\bar B\sim\bigoplus_v B_v$ be unit-weight
    cyclic sector decompositions whose sectors are left-canonical. Let the
    specified cyclic projectors be $P_u$ and $Q_v$, and let the specified
    corner identifications realize the blocked sector tensors by
    $A_u^\alpha=P_u\bar A^\alpha P_u$ and
    $B_v^\alpha=Q_v\bar B^\alpha Q_v$.
    Assume
    $d_u\ne0$ and $A_u$ is normal for every $u$. Fix $q\in\mathbb Z_m$ and
    assume that for every $u$ there are $d_u=e_{u+q}$, $\zeta_u\ne0$, and
    $X_u\in\GL_{d_u}(\C)$ such that
    $B_{u+q}^\alpha=\zeta_u X_u A_u^\alpha X_u^{-1}$ for every blocked
    physical index $\alpha$. The source obtains a phase-only product identity
    after absorbing these sector gauges into the $B$-sectors: with $Q_v$ the
    cyclic sector projectors of $B$, it writes
    $\widetilde B_v^i
        =U_v Q_v B^i Q_{v+1} U_{v+1}^{\dagger}$.
    The contraction input is then, for every physical block
    $(i_1,\ldots,i_m)$,
    \begin{align}
        A_u^{i_1}A_{u+1}^{i_2}\cdots A_{u+m-1}^{i_m}
         & =
        e^{i\lambda_{u+q}}
        \widetilde B_{u+q}^{i_1}\widetilde B_{u+q+1}^{i_2}
        \cdots\widetilde B_{u+q+m-1}^{i_m}.
        \notag
    \end{align}
    % arXiv:1708.00029 Appendix A: eq:Fu, eq:Omegauprop, eq:resultprop.
    Choose $N_0$ so that the $N_0$-fold repetition of each cyclic
    $m$-site sector product is injective. For
    $\mathbf i=(i_1,\ldots,i_{mN_0})$, form the repeated product
    \begin{align}
        F_u^{\mathbf i}
         & =
        (A_u^{i_1}\cdots A_{u+m-1}^{i_m})
        \cdots
        (A_u^{i_{m(N_0-1)+1}}\cdots A_{u+m-1}^{i_{mN_0}}),
        \notag
    \end{align}
    and choose right inverses $\Omega_u$ satisfying
    \begin{align}
        \sum_{\mathbf i}(\Omega_u^{\mathbf i})_{\alpha,\beta}
        (F_u^{\mathbf i})_{\alpha',\beta'}
         & =
        \delta_{\alpha,\alpha'}\delta_{\beta,\beta'}.
        \notag
    \end{align}
    Let $\widetilde F_{u+q+r}^{\mathbf i}$ denote the same cyclic repeated
    product formed from
    $\widetilde B_{u+q+r},\ldots,\widetilde B_{u+q+r+m-1}$. For every
    $r\in\mathbb Z_m$, the block identity gives
    \begin{align}
        F_{u+r}^{\mathbf i}
         & =
        e^{iN_0\lambda_{u+q+r}}\widetilde F_{u+q+r}^{\mathbf i},
        \notag \\
        \widetilde F_{u+q+r}^{\mathbf i}
         & =
        e^{-iN_0\lambda_{u+q+r}}F_{u+r}^{\mathbf i}.
        \notag
    \end{align}
    % Local fix (docs/paper-gaps/dccsp17_periodic_overlap_route_alignment.tex):
    % arXiv:1708.00029, lines 1041--1060, displays $m-1$ inverses although the
    % repetition count and phase formula require one inverse for each inserted product.
    Let $\mathbf i_1,\ldots,\mathbf i_m$ be independent words of length $mN_0$,
    one for each inserted repeated product. Sum each $\mathbf i_r$ in the
    contraction by the right-inverse relation for $\Omega_{u+r}$:
    \begin{align}
         & A_u^{i_1}F_{u+1}^{\mathbf i_1}
        A_{u+1}^{i_2}F_{u+2}^{\mathbf i_2}
        \cdots
        A_{u+m-1}^{i_m}F_u^{\mathbf i_m}
        \notag                                                                \\
         & \quad\xrightarrow{\ \Omega_{u+1}\otimes\cdots\otimes\Omega_{u+m-1}
            \otimes\Omega_u\ }
        A_u^{i_1}\otimes A_{u+1}^{i_2}\otimes\cdots\otimes A_{u+m-1}^{i_m}.
        \notag
    \end{align}
    Expand the words $\mathbf i_1,\ldots,\mathbf i_m$, and write $j_{s,t}$ for
    the physical index in the $t$th sector of the $s$th consecutive $m$-site
    block of the expanded word, where $0\leq s\leq mN_0$ and $1\leq t\leq m$.
    Since each inserted product $F_{u+r}^{\mathbf i_r}$ has length $mN_0$,
    insertion shifts the cyclic sector position by $mN_0\equiv 0\pmod m$.
    Hence every consecutive $m$-site block begins at the sector $u$, and the
    interleaved word is
    \begin{align}
         & A_u^{i_1}F_{u+1}^{\mathbf i_1}
        A_{u+1}^{i_2}F_{u+2}^{\mathbf i_2}
        \cdots
        A_{u+m-1}^{i_m}F_u^{\mathbf i_m}
        \notag                            \\
         & \quad=
        (A_u^{j_{0,1}}\cdots A_{u+m-1}^{j_{0,m}})
        \cdots
        (A_u^{j_{mN_0,1}}\cdots A_{u+m-1}^{j_{mN_0,m}}).
        \notag
    \end{align}
    Applying the block identity at the sector $u$ to each of these $m$-site
    blocks gives
    \begin{align}
         & A_u^{i_1}F_{u+1}^{\mathbf i_1}
        A_{u+1}^{i_2}F_{u+2}^{\mathbf i_2}
        \cdots
        A_{u+m-1}^{i_m}F_u^{\mathbf i_m}
        \notag                            \\
         & \quad=
        e^{i\lambda_{u+q}(mN_0+1)}
        \widetilde B_{u+q}^{i_1}\widetilde F_{u+q+1}^{\mathbf i_1}
        \widetilde B_{u+q+1}^{i_2}\widetilde F_{u+q+2}^{\mathbf i_2}
        \cdots
        \widetilde B_{u+q+m-1}^{i_m}\widetilde F_{u+q}^{\mathbf i_m}.
        \label{eq:pft_intermediate_block_identity}
    \end{align}
    For each $r\in\mathbb Z_m$, substituting $\mathbf i=\mathbf i_r$ in
    (\ref{eq:pft_intermediate_block_identity}) gives
    $\widetilde F_{u+q+r}^{\mathbf i_r}
        =e^{-iN_0\lambda_{u+q+r}}F_{u+r}^{\mathbf i_r}$.
    Therefore the right-inverse relation gives the per-sector contraction
    identity
    \begin{align}
        \sum_{\mathbf i_r}
        (\Omega_{u+r}^{\mathbf i_r})_{\alpha,\beta}
        (\widetilde F_{u+q+r}^{\mathbf i_r})_{\alpha',\beta'}
         & =
        e^{-iN_0\lambda_{u+q+r}}
        \delta_{\alpha,\alpha'}\delta_{\beta,\beta'}.
        \notag
    \end{align}
    Since $r\mapsto u+q+r$ is a permutation of $\mathbb Z_m$,
    $\sum_{r\in\mathbb Z_m}\lambda_{u+q+r}
        =\sum_{r\in\mathbb Z_m}\lambda_r$.
    Applying
    $\Omega_{u+1}\otimes\cdots\otimes\Omega_{u+m-1}\otimes\Omega_u$ to both
    sides of (\ref{eq:pft_intermediate_block_identity}) and collecting the
    sector phases gives
    \begin{align}
        A_u^{i_1}\otimes A_{u+1}^{i_2}\otimes\cdots\otimes A_{u+m-1}^{i_m}
         & =
        e^{i\lambda_{u+q}(mN_0+1)-iN_0\sum_{r\in\mathbb Z_m}\lambda_r}
        \widetilde B_{u+q}^{i_1}\otimes\widetilde B_{u+q+1}^{i_2}
        \otimes\cdots\otimes\widetilde B_{u+q+m-1}^{i_m}
        \notag \\
         & =
        e^{i\eta_{u+q}}
        \widetilde B_{u+q}^{i_1}\otimes\widetilde B_{u+q+1}^{i_2}
        \otimes\cdots\otimes\widetilde B_{u+q+m-1}^{i_m}.
        \notag
    \end{align}
    Here
    $\eta_{u+q}=\lambda_{u+q}(mN_0+1)
        -N_0\sum_{r\in\mathbb Z_m}\lambda_r$.
    Then there are a scalar $\xi$ with $|\xi|=1$ and a unitary matrix $U$
    such that
    $A^i=\xi UB^iU^\dagger$ for every physical index $i$.
\end{lemma}

\begin{theorem}[Sector match yields repeated-block equivalence]%
    \label{thm:periodic_overlap_gauge_equiv_sector_match}
    \lean{MPSTensor.periodicOverlap_gaugeEquiv_of_sector_match}
    \leanok
    \uses{lem:sector_match_propagation, lem:sector_tensor_proportional_blocked_match, def:cyclic_sector_decomp}
    Let $A$ and $B$ be periodic tensors with the same period $m$. Let
    $\bar A\sim\bigoplus_u A_u$ and $\bar B\sim\bigoplus_v B_v$ be unit-weight
    cyclic sector decompositions whose sectors are left-canonical, with
    specified cyclic projectors $P_u,Q_v$ and corner identifications satisfying
    $A_u^\alpha=P_u\bar A^\alpha P_u$ and
    $B_v^\alpha=Q_v\bar B^\alpha Q_v$.
    Assume
    every sector $A_u$ has nonzero virtual dimension. If there are sectors
    $u_0,v_0$ with $d_{u_0}=e_{v_0}$ and
    $A_{u_0}\sim_{\mathrm{gp}}B_{v_0}$ after identifying their virtual spaces,
    then $A$ and $B$ are repeated blocks.
\end{theorem}

\begin{theorem}[Different bond dimensions imply asymptotic orthogonality]%
    \label{thm:periodic_overlap_zero_ne_dim}
    \lean{MPSTensor.periodicOverlap_tendsto_zero_of_ne_dim}
    \leanok
    \uses{def:periodic_tensor}
    If two periodic tensors have different bond dimensions, then
    $\braket{V^{(N)}(A)}{V^{(N)}(B)} \to 0$ as $N\to\infty$.
\end{theorem}

\begin{theorem}[Periodic overlap dichotomy]\label{thm:periodic_overlap_dichotomy}
    \lean{MPSTensor.periodicOverlapDichotomy}
    \leanok
    \uses{thm:periodic_overlap_zero_ne_period, thm:periodic_overlap_zero_no_sector_match, thm:periodic_overlap_zero_ne_dim, thm:periodic_overlap_gauge_equiv_sector_match}
    Let $A$ and $B$ be periodic tensors. Then at least one of the following
    alternatives holds:
    \begin{enumerate}
        \item $\braket{V^{(N)}(A)}{V^{(N)}(B)} \to 0$ as $N\to\infty$;
        \item $A$ and $B$ are related by the repeated-block relation, i.e. they
              differ by a unit-modulus phase and an invertible gauge.
    \end{enumerate}
\end{theorem}
\begin{proof}\leanok
    \uses{thm:periodic_overlap_zero_ne_period, thm:periodic_overlap_zero_no_sector_match, thm:periodic_overlap_zero_ne_dim, thm:periodic_overlap_gauge_equiv_sector_match}
    If $m_A\ne m_B$, then
    $\lim_{N\to\infty}O_{AB}(N)=0$ by
    Theorem~\ref{thm:periodic_overlap_zero_ne_period}. Assume $m_A=m_B=m$.
    If $D_A\ne D_B$, then
    $\lim_{N\to\infty}O_{AB}(N)=0$ by
    Theorem~\ref{thm:periodic_overlap_zero_ne_dim}. Assume $D_A=D_B$.

    Choose period-blocked sector decompositions
    \begin{align}
        \bar A & \sim\bigoplus_{u\in\mathbb Z_m}A_u,
        \notag                                       \\
        \bar B & \sim\bigoplus_{v\in\mathbb Z_m}B_v.
        \notag
    \end{align}
    If no pair $(u,v)$ satisfies
    $d_u=e_v$ and $A_u\sim_{\mathrm{gp}}B_v$, then
    Theorem~\ref{thm:periodic_overlap_zero_no_sector_match} gives
    $\lim_{N\to\infty}O_{AB}(N)=0$.

    Otherwise there are $u_0,v_0\in\mathbb Z_m$ such that
    $d_{u_0}=e_{v_0}$ and
    $A_{u_0}\sim_{\mathrm{gp}}B_{v_0}$.
    Theorem~\ref{thm:periodic_overlap_gauge_equiv_sector_match} then gives
    repeated-block equivalence between $A$ and $B$.
\end{proof}

\begin{theorem}[Gram criterion for a periodic family]
    \label{thm:periodic_family_li_of_cross_overlap_decay}
    \lean{MPSTensor.%
        periodicFamily_eventuallyLinearlyIndependent_of_cross_overlap_tendsto_zero}
    \leanok
    \uses{def:periodic_tensor}
    Let $\{A_j\}_{j\in J}$ be a finite family of periodic tensors whose
    periods divide a common positive integer $p$. If
    \begin{align}
        \braket{V^{(N)}(A_i)}{V^{(N)}(A_j)}\longrightarrow0
        \qquad (i\ne j),
        \notag
    \end{align}
    then the vectors $\{\ket{V^{(pN)}(A_j)}\}_{j\in J}$ are linearly
    independent for every sufficiently large $N$.
\end{theorem}
\begin{proof}\leanok
    \uses{thm:periodic_bond_dimension_pos, thm:periodic_self_overlap_tendsto}
    The Gram matrix converges to the diagonal matrix whose $j$-th diagonal
    entry is the positive period of $A_j$. Its determinant therefore has a
    nonzero limit, so the Gram matrix is invertible for every sufficiently
    large $N$.
\end{proof}

\begin{theorem}[Eventual linear independence of periodic vectors]%
    \label{thm:periodic_basis_eventually_li}
    \lean{MPSTensor.periodicBasis_eventuallyLinearlyIndependent}
    \leanok
    \uses{def:periodic_tensor}
    For a finite family $\{A_j\}_{j\in J}$ of pairwise non-repeated
    periodic tensors, in the sense that no two are related by a gauge
    transformation and a unit-modulus scalar, whose periods divide a common
    positive integer $p$,
    there is $N_0$ such that for every $N \ge N_0$, the map
    \begin{align}
        \Phi_N\colon\C^J & \longrightarrow(\C^d)^{\otimes pN},
        \notag                                                 \\
        (c_j)_{j\in J}   & \longmapsto
        \sum_{j\in J}c_j\ket{V^{(pN)}(A_j)}
        \notag
    \end{align}
    satisfies $\ker\Phi_N=\{0\}$.

    \emph{Scope restriction.}  The source consequence following the
    periodic-overlap dichotomy
    \cite[lines~589--609]{DeLasCuevas2017Irreducible} also states that, for
    an ambient tensor $A$ in irreducible form having $\{A_j\}_{j\in J}$ as
    its basis of periodic tensors, the non-zero vectors $\ket{V^{(N)}(A_j)}$
    span the vector $\ket{V^{(N)}(A)}$;
    this spanning assertion is what the paper cites as justification for
    the phrase ``basis of periodic vectors''. The theorem recorded here
    omits the spanning assertion. The corresponding paper-gap
    note, \cite{gap:dccsp17_periodic_overlap_route_alignment},
    Section ``Scope restriction: independence without spanning'', records
    this restriction and its elimination plan.
\end{theorem}

\begin{proof}\leanok
    \uses{thm:periodic_bond_dimension_pos, thm:periodic_family_li_of_cross_overlap_decay, thm:periodic_overlap_dichotomy}
    Periodicity makes every bond dimension positive. For distinct $i,j$, the
    repeated-block alternative in
    Theorem~\ref{thm:periodic_overlap_dichotomy} contradicts pairwise
    non-repetition, so the cross overlap tends to zero.
    Theorem~\ref{thm:periodic_family_li_of_cross_overlap_decay} now gives the
    claimed independence along the common-period subsequence.
\end{proof}

\begin{theorem}[Independence of the non-zero periodic vectors]%
    \label{thm:periodic_basis_nonzero_subfamily_li}
    \lean{MPSTensor.periodicBasis_nonzero_subfamily_eventuallyLinearlyIndependent}
    \leanok
    \uses{def:periodic_tensor}
    Let $\{A_j\}_{j\in J}$ be a finite family of pairwise non-repeated
    periodic tensors, of periods $m_j$. There is an $N_0$ such that for every
    single length $N\ge N_0$ the vectors
    \begin{align}
        \{\ket{V^{(N)}(A_j)} : j\in J,\ m_j\mid N\}
        \notag
    \end{align}
    are linearly independent. Since $\ket{V^{(N)}(A_j)}$ vanishes unless
    $m_j\mid N$, this is the assertion of
    \cite[lines~604--608]{DeLasCuevas2017Irreducible} that beyond some length
    the non-zero members of the family are independent, in the form used in
    the proof of the equal case at
    \cite[lines~677--679]{DeLasCuevas2017Irreducible}.
\end{theorem}

\begin{proof}\leanok
    \uses{thm:periodic_bond_dimension_pos, thm:periodic_family_li_of_cross_overlap_decay, thm:periodic_overlap_dichotomy}
    There are finitely many subfamilies. For each subfamily $S\subseteq J$
    apply Theorem~\ref{thm:periodic_family_li_of_cross_overlap_decay} at the
    least common multiple $p_S$ of the periods occurring in $S$, the
    cross-overlap decay coming as before from
    Theorem~\ref{thm:periodic_overlap_dichotomy} and pairwise non-repetition;
    this gives a threshold $M_S$. Let $L$ be the least common multiple of all
    the periods and take $N_0$ to be the largest of the numbers $L\,M_S$.
    Given $N\ge N_0$, let $S$ be the set of indices whose period divides $N$.
    Every period occurring in $S$ divides $N$, so $p_S\mid N$; writing
    $N=p_Sn$ and using $p_S\le L$ gives $L\,M_S\le N\le Ln$, hence $n\ge M_S$
    and the subfamily is independent at the length $N=p_Sn$.
\end{proof}

\begin{remark}[Comparison with the blocking approach]\notready
    \label{rem:periodic_vs_blocking_ft}
    Theorem~\ref{thm:periodic_ft} and the equal-MPV corollary
    of Section~\ref{sec:ft_equal_proportional_cases} both characterize the gauge freedom
    between two tensors with equal MPV families, but they do so at
    different levels:

    \begin{enumerate}
        \item \emph{Blocking approach} (Section~\ref{sec:ft_equal_proportional_cases},
              following \cite{Cirac2016MPDO_arXiv}):
              for a common multiple $p$ of the periods,
              \begin{align}
                  A^{[p]} & \sim\bigoplus_a C_a,
                  \notag                         \\
                  B^{[p]} & \sim\bigoplus_b D_b,
                  \notag
              \end{align}
              with all $C_a,D_b$ normal. The Fundamental Theorem gives, after a
              permutation of the normal summands,
              \begin{align}
                  D_{\pi(a)}^\alpha
                   & =
                  X_a C_a^\alpha X_a^{-1}.
                  \notag
              \end{align}
              This is a gauge statement for $A^{[p]}$ and $B^{[p]}$; the cyclic
              phases have been absorbed into the period-$p$ words.

        \item \emph{Periodic approach} (Theorem~\ref{thm:periodic_ft},
              following \cite{DeLasCuevas2017Irreducible}):
              keep the cyclic sectors
              \begin{align}
                  A^i & =\sum_{u\in\mathbb Z_m}P_uA^iP_{u+1},
                  \notag                                      \\
                  B^i & =\sum_{v\in\mathbb Z_m}Q_vB^iQ_{v+1}.
                  \notag
              \end{align}
              Sector matching gives phases $\phi_v$, a shift $q$, and unitaries
              $U_v$ such that
              \begin{align}
                  P_uA^iP_{u+1}
                   & =
                  e^{i\xi}e^{i\phi_{u+q}}
                  U_{u+q}Q_{u+q}B^iQ_{u+q+1}U_{u+q+1}^{\dagger}
                  e^{-i\phi_{u+q+1}}.
                  \notag
              \end{align}
              In the equal-MPV theorem the proportional phase is absorbed, so
              one obtains $ZA^i=UB^iU^{-1}$ with $Z^m=\Id$; the trace only sees
              powers of the corresponding $m$-th roots of unity.
    \end{enumerate}

    Chapters~\ref{ch:canonical} and~\ref{ch:ft_proof} follow only the
    blocking argument: remove periodicity by blocking, then apply the
    Fundamental Theorem. The periodic extension treats periodic blocks directly.
\end{remark}

\begin{theorem}[Peripheral proportional case from positive-length MPV equality]
    \label{thm:peripheral_periodic_ft_proportional_same}
    \lean{MPSTensor.peripheralProportionalCase_periodicFT_of_sameMPV₂Pos}
    \leanok
    \uses{def:periodic_tensor, def:same_mpv2_pos, thm:periodic_overlap_dichotomy, thm:periodic_self_overlap_tendsto}
    Let $A$ and $B$ be periodic tensors with the same matrix product vectors
    at every positive length.
    Then their bond dimensions agree, and after identifying the bond spaces,
    $A$ and $B$ are repeated blocks.
\end{theorem}
\begin{proof}\leanok
    By Theorem~\ref{thm:periodic_overlap_dichotomy}, either the overlap
    $\braket{V^{(N)}(A)}{V^{(N)}(B)}$ tends to $0$ or the bond dimensions agree
    and $A$ and $B$ are repeated blocks. In the decay case,
    equality of the MPV families gives, for every $N>0$,
    \begin{align}
        \braket{V^{(N)}(A)}{V^{(N)}(B)}
         & =
        \braket{V^{(N)}(A)}{V^{(N)}(A)}.
        \label{eq:pft_decay_mpv_equality}
    \end{align}
    Along multiples of the period $m_a$, (\ref{eq:pft_decay_mpv_equality})
    would imply
    \begin{align}
        0
         & =
        \lim_{N\to\infty}
        \braket{V^{(m_aN)}(A)}{V^{(m_aN)}(B)}
        =
        \lim_{N\to\infty}
        \braket{V^{(m_aN)}(A)}{V^{(m_aN)}(A)}
        =m_a,
        \notag
    \end{align}
    contradicting $m_a>0$.
\end{proof}

\begin{lemma}[A pure similarity is a repeated-block relation]
    \label{lem:pure_gauge_het_repeated_blocks}
    \lean{MPSTensor.GaugeEquiv.toHetRepeatedBlocks}
    \leanok
    If $B^i=XA^iX^{-1}$ for an invertible matrix $X$, then $A$ and $B$ are
    repeated blocks with equal bond dimension and unit phase.
\end{lemma}
\begin{proof}\leanok
    Invert the pure gauge and take the repeated-block phase to be one.
\end{proof}

\begin{theorem}[Repeated blocks scale their vectors by a phase power]
    \label{thm:repeated_blocks_phase_power_mpv}
    \lean{MPSTensor.HetRepeatedBlocks.exists_unit_phase_power_mpv}
    \leanok
    Let $A$ and $B$ be repeated blocks.  Then there is a scalar $\zeta$ with
    $|\zeta|=1$ such that, for every length $N$,
    \begin{align}
        \ket{V^{(N)}(A)}=\zeta^{N}\ket{V^{(N)}(B)}.
        \notag
    \end{align}
    The proportionality scalar therefore depends on the length.  Already the
    one-dimensional blocks with a single physical index, $B^0=1$ and
    $A^0=\zeta$, satisfy $\ket{V^{(N)}(A)}=\zeta^{N}\ket{V^{(N)}(B)}$, so no
    scalar independent of the length works unless $\zeta=1$.
\end{theorem}
\begin{proof}
    \uses{thm:gauge_phase_mpv_scaling}
    \leanok
    Write $A^i=\zeta YB^iY^{-1}$ with $|\zeta|=1$.  A closed chain of length
    $N$ is the trace of a product of $N$ such factors.  The neighbouring
    similarities cancel and the trace is invariant under the outer
    conjugation, while each site contributes one factor $\zeta$.
\end{proof}

\begin{theorem}[Absorbing the phase of a repeated-block relation]
    \label{thm:repeated_blocks_phase_absorption}
    \lean{MPSTensor.HetRepeatedBlocks.exists_phase_rescaling_sameMPV₂Pos}
    \leanok
    Let $A$ and $B$ be repeated blocks.  Then there is a scalar $\xi$ with
    $|\xi|=1$ such that $\xi A$ and $B$ generate the same vector at every
    positive length,
    \begin{align}
        \ket{V^{(N)}(\xi A)}=\ket{V^{(N)}(B)}.
        \notag
    \end{align}
    The phase belongs to the matched pair $A,B$; distinct matched pairs of two
    bases of periodic tensors may require distinct phases, so no single phase
    is asserted for the assembled tensors
    \cite[lines~286--305, 625--627]{DeLasCuevas2017Irreducible}.
\end{theorem}
\begin{proof}
    \uses{thm:repeated_blocks_phase_power_mpv, thm:gauge_phase_mpv_scaling}
    \leanok
    Take $\xi$ to be the reciprocal of the phase of the preceding theorem, which
    is nonzero because it has modulus one.  Rescaling every matrix of a block
    by $\xi$ multiplies its length-$N$ vector by $\xi^{N}$, and the two powers
    cancel.
\end{proof}

\begin{theorem}[Unit-norm scaling preserves the transfer map]
    \label{thm:transfer_map_smul_unit_norm}
    \lean{MPSTensor.transferMap_smul_eq_of_norm_eq_one}
    \leanok
    If $|c|=1$, then the tensors $A$ and $cA$ have the same transfer map.
\end{theorem}
\begin{proof}
    \uses{lem:transfer_map_smul}
    \leanok
    Scaling every tensor letter by $c$ multiplies the transfer map by
    $c\overline c=|c|^2=1$.
\end{proof}

\begin{theorem}[Repeated blocks share a peripheral spectrum]
    \label{thm:repeated_blocks_peripheral_spectrum}
    \lean{MPSTensor.RepeatedBlocks.peripheralEigenvalues_transferMap_eq}
    \leanok
    If $A$ and $B$ are repeated blocks, then their transfer maps have the same
    unit-circle eigenvalues.
\end{theorem}
\begin{proof}
    \uses{thm:transfer_map_smul_unit_norm, thm:cpsv_gauge_transfer_similarity, %
        lem:peripheral_spectrum_similarity}
    \leanok
    Write $A^i=\xi YB^iY^{-1}$ with $|\xi|=1$.  The transfer map of $\xi C$ is
    $|\xi|^2$ times the transfer map of $C$, hence equals it, and the transfer
    map of $YB^iY^{-1}$ is the congruence conjugate
    $Y\mathcal{E}_B(Y^{-1}\cdot Y^{-\dagger})Y^{\dagger}$ of the transfer map
    of $B$.  Conjugation by an invertible linear map leaves the spectrum
    unchanged.
\end{proof}

\begin{theorem}[A root-of-unity set determines its positive exponent]
    \label{thm:roots_of_unity_determine_period}
    \lean{MPSTensor.period_eq_of_setOf_pow_eq_one}
    \leanok
    Let $m,n>0$. If
    $\{z\in\C:z^m=1\}=\{z\in\C:z^n=1\}$, then $m=n$.
\end{theorem}
\begin{proof}\leanok
    A primitive $m$-th root belongs to the second set, so $m$ divides $n$.
    Interchanging $m$ and $n$ gives the reverse divisibility.
\end{proof}

\begin{theorem}[A common peripheral spectrum determines the period]
    \label{thm:spectral_period_eq_of_peripheral_spectrum}
    \lean{MPSTensor.IsSpectrallyPeriodic.period_eq_of_peripheralEigenvalues_eq}
    \leanok
    \uses{def:spectral_periodic_tensor}
    Let $A$ and $B$ be spectrally periodic tensors of periods $m$ and $n$. If
    their transfer maps have the same peripheral spectrum, then $m=n$.
\end{theorem}
\begin{proof}
    \uses{thm:roots_of_unity_determine_period}
    \leanok
    The two spectra are respectively the sets of $m$-th and $n$-th roots of
    unity, so the preceding theorem applies.
\end{proof}

\begin{theorem}[Repeated periodic blocks have the same period]
    \label{thm:repeated_blocks_period_eq}
    \lean{MPSTensor.IsSpectrallyPeriodic.period_eq_of_hetRepeatedBlocks}
    \leanok
    \uses{def:spectral_periodic_tensor}
    Let $A$ and $B$ be spectrally periodic tensors of periods $m$ and $n$.  If
    $A$ and $B$ are repeated blocks, then $m=n$.
\end{theorem}
\begin{proof}
    \uses{thm:repeated_blocks_peripheral_spectrum, %
        thm:spectral_period_eq_of_peripheral_spectrum}
    \leanok
    Repeated blocks have the same peripheral spectrum. The preceding theorem
    then identifies their periods.
\end{proof}

\begin{lemma}[Pure similarities preserve periodic overlap hypotheses]
    \label{lem:periodic_overlap_hypothesis_gauge_transport}
    \lean{MPSTensor.PeriodicOverlapHypothesis.of_gaugeEquiv}
    \leanok
    Let $\{A_j\}_{j\in J}$ and $\{B_k\}_{k\in K}$ be finite tensor
    families, and suppose that $A_j'$ is purely similar to $A_j$ and $B_k'$ is
    purely similar to $B_k$ for every $j,k$.  Assume that every tensor in each
    primed family has a non-decaying mixed-overlap partner in the other primed
    family, and that every non-decaying primed pair is a repeated pair.  Then
    the same three assertions hold for the two original families.
\end{lemma}
\begin{proof}
    \uses{lem:pure_gauge_het_repeated_blocks, thm:gauge_invariance}
    \leanok
    Cyclicity of the trace makes every closed-chain coefficient, and hence
    every mixed overlap, invariant under the two similarities.  Thus decay and
    non-decay are unchanged.  For a non-decaying pair, compose its primed
    repeated-block relation with the two pure similarities.
\end{proof}

\begin{theorem}[Non-decaying mixed overlaps for normalized bases of periodic
        tensors]
    \label{thm:periodic_multiplicity_nondecaying_overlap}
    \lean{MPSTensor.PeriodicOverlapHypothesis.ofSectorDecompositions}
    \leanok
    \uses{def:nonzero_proportional_mpv}
    Let
    \begin{align}
        A^i & =\bigoplus_{j\in J}(R_j\otimes A_j^i),
        \notag                                       \\
        B^i & =\bigoplus_{k\in K}(S_k\otimes B_k^i),
        \notag
    \end{align}
    where $R_j=\operatorname{diag}(\mu_{j,1},\ldots,\mu_{j,r_j})$ and
    $S_k=\operatorname{diag}(\nu_{k,1},\ldots,\nu_{k,s_k})$ have nonzero
    complex diagonal entries.  Suppose that the finite families
    $\{A_j\}_{j\in J}$ and $\{B_k\}_{k\in K}$ are pairwise non-repeated
    families of left-canonical periodic tensors.  If, for every positive
    integer $N$, there is a scalar $c_N\ne0$ such that
    \begin{align}
        \ket{V^{(N)}(A)}=c_N\ket{V^{(N)}(B)}.
        \notag
    \end{align}
    Then every $A_j$ has an index $k\in K$ for which
    \begin{align}
        \braket{V^{(N)}(A_j)}{V^{(N)}(B_k)}
        \not\longrightarrow 0,
        \notag
    \end{align}
    and, conversely, every $B_k$ has such a partner in the first family.
    Moreover, whenever the overlap between $A_j$ and $B_k$ does not tend to
    zero, the two blocks have equal bond dimension and there are an invertible
    matrix $Y$ and a scalar $\zeta$ with $|\zeta|=1$ such that, for every
    physical index $i$,
    \begin{align}
        A_j^i & =\zeta YB_k^iY^{-1}.
        \notag
    \end{align}

    \emph{Normalization boundary.} The source initially assumes that the transfer
    maps of the basis tensors are irreducible and have spectral radius one,
    and then obtains the trace-preserving normalization by a block-diagonal
    similarity
    \cite[lines~313--332]{DeLasCuevas2017Irreducible}.  The present theorem
    starts with the normalized basis tensors.  The multiplicity matrices are
    otherwise arbitrary diagonal matrices with nonzero complex entries, the
    displayed direct sums are literal, and proportionality is assumed only at
    positive lengths.  Theorem~\ref{thm:spectral_periodic_normalization}
    constructs each normalizing similarity, while
    Theorem~\ref{thm:spectral_periodic_sector_normalization} assembles them
    without changing the multiplicity matrices.  Theorem~\ref{thm:%
        periodic_spectral_multiplicity_nondecaying_overlap} below applies the
    present result and returns its conclusion to the unnormalized blocks.
\end{theorem}
\begin{proof}\leanok
    \uses{lem:finite_coefficient_isolation, lem:nonzero_proportional_mpv_basic, thm:gauge_phase_mpv_scaling, thm:mpv_decomposition, thm:periodic_bond_dimension_pos, thm:periodic_family_li_of_cross_overlap_decay, thm:periodic_overlap_dichotomy, lem:sector_decomposition_coeff_not_eventually_zero}
    Fix $j_0\in J$ and suppose that its overlap with every $B_k$ tends to
    zero. Let $T\subseteq J$ be the set of all indices with this property.
    For every $j\notin T$, choose $k(j)$ for which the mixed overlap does not
    decay. The periodic overlap dichotomy identifies $A_j$ with $B_{k(j)}$
    up to a gauge transformation and a phase; consequently their length-$N$
    vectors are scalar multiples of one another.

    Choose a common positive multiple $p$ of all periods. Along the lengths
    $pN$, define the coproduct-indexed family $w_N\colon T\amalg K\to
        (\mathbb C^d)^{\otimes pN}$ by
    $w_N(\operatorname{inl}(j))=\ket{V^{(pN)}(A_j)}$ and
    $w_N(\operatorname{inr}(k))=\ket{V^{(pN)}(B_k)}$. This family is eventually
    linearly independent. Indeed, distinct vectors within either summand have
    decaying overlaps by non-repetition and the overlap dichotomy, while the
    mixed overlaps vanish by the definition of $T$.

    Expand the proportionality relation at length $pN$ using the grouped
    coefficients
    \begin{align}
        \operatorname{tr}(R_j^{pN}) & =\sum_q(\mu_{j,q}^{p})^N, &
        \operatorname{tr}(S_k^{pN}) & =\sum_r(\nu_{k,r}^{p})^N.
        \notag
    \end{align}
    Replace each term indexed by $j\notin T$ by its scalar multiple of the
    corresponding $B_{k(j)}$ term, and collect terms with the same value of
    $k(j)$. Linear independence forces
    $\operatorname{tr}(R_{j_0}^{pN})=0$ for every sufficiently large $N$.
    Since each $\mu_{j_0,q}^p$ is nonzero, geometric extrapolation gives the
    same identity at exponent zero.  This says $r_{j_0}=0$, contradicting the
    positive size of $R_{j_0}$. Hence $j_0$ has a non-decaying partner.
    Interchanging $A$ and $B$, inverting the nonzero scalars $c_N$, and using
    conjugate symmetry of the mixed overlap gives the converse assertion. For
    any non-decaying pair, the periodic overlap dichotomy excludes its decay
    alternative and supplies the repeated-block relation.
\end{proof}

\begin{theorem}[Non-decaying mixed overlaps for spectrally periodic bases]
    \label{thm:periodic_spectral_multiplicity_nondecaying_overlap}
    \lean{MPSTensor.PeriodicOverlapHypothesis.%
        ofSpectrallyPeriodicSectorDecompositions}
    \leanok
    \uses{def:nonzero_proportional_mpv, def:spectral_periodic_tensor}
    Let
    \begin{align}
        A^i & =\bigoplus_{j\in J}(R_j\otimes A_j^i), &
        B^i & =\bigoplus_{k\in K}(S_k\otimes B_k^i),
        \notag
    \end{align}
    where the diagonal entries of every $R_j$ and $S_k$ are nonzero.  Suppose
    that the two representative families are pairwise non-repeated.  Assign
    positive integers $m_j$ and $n_k$ to their representatives, and suppose
    that every transfer map is irreducible, has spectral radius one, and has
    peripheral spectrum respectively
    $\{z\in\C:z^{m_j}=1\}$ or $\{z\in\C:z^{n_k}=1\}$.  If, for every positive
    $N$, there is a scalar $c_N\ne0$ such that
    \begin{align}
        \ket{V^{(N)}(A)} & =c_N\ket{V^{(N)}(B)},
        \notag
    \end{align}
    then every representative in either family has a representative in the
    other family whose mixed overlap does not tend to zero.  Every such pair
    has equal bond dimension and differs by an invertible similarity and a
    unit-modulus scalar.
\end{theorem}
\begin{proof}
    \uses{lem:periodic_overlap_hypothesis_gauge_transport, lem:pure_gauge_het_repeated_blocks, thm:gauge_invariance, thm:periodic_multiplicity_nondecaying_overlap, thm:spectral_periodic_sector_normalization}
    \leanok
    Normalize all representatives by their pure Perron gauges and assemble
    the gauges block diagonally.  The multiplicity matrices and their grouped
    coefficients remain unchanged.  Pure similarity preserves the total
    positive-length proportionality and pairwise non-repetition, so
    Theorem~\ref{thm:periodic_multiplicity_nondecaying_overlap} applies to the
    normalized families.  Finally apply
    Lemma~\ref{lem:periodic_overlap_hypothesis_gauge_transport}.
\end{proof}

\begin{theorem}[Scalar-weight specialization of the mixed-overlap theorem]
    \label{thm:periodic_scalar_weight_nondecaying_overlap}
    \lean{MPSTensor.PeriodicOverlapHypothesis.ofIsIrreducibleForm}
    \leanok
    \uses{def:irreducible_form, def:nonzero_proportional_mpv}
    Let $A$ and $B$ each have the same length-$N$ matrix-product vector, for
    every $N\geq0$, as a scalar-weight direct sum formed from a finite,
    pairwise non-repeated family of left-canonical periodic tensors
    $\{A_j\}_{j\in J}$ and $\{B_k\}_{k\in K}$, with a positive real scalar
    weight attached to each block. Suppose that, for every positive integer
    $N$, there is a scalar $c_N\ne0$ such that
    \begin{align}
        \ket{V^{(N)}(A)}=c_N\ket{V^{(N)}(B)}.
        \notag
    \end{align}
    Then every block in either family has a block in the other family whose
    mixed overlap does not tend to zero.  Every such non-decaying pair has
    equal bond dimension and differs by an invertible similarity and a
    unit-modulus scalar.

    \emph{Scope restriction.} This formulation records one positive scalar
    per listed block and equality with the corresponding direct sum also at
    the empty word.  Thus it is a specialization of the preceding
    multiplicity-bearing theorem and is not itself the unrestricted
    irreducible-form statement of
    \cite[Theorem~3.4]{DeLasCuevas2017Irreducible}.
\end{theorem}
\begin{proof}\leanok
    Repeat the coefficient argument above with the singleton coefficients
    $\mu_j^N$ and $\nu_k^N$.  Their nonvanishing at every length replaces the
    geometric-extrapolation step.  Equality with the two scalar-weight direct
    sums transports the positive-length proportionality relation to the
    corresponding block-diagonal tensors.
\end{proof}

\begin{theorem}[Finite matching from the periodic overlap dichotomy]
    \label{thm:periodic_conditional_block_matching}
    \lean{MPSTensor.fundamentalTheorem_periodic_proportional}
    \leanok
    Let $\{A_j\}_{j\in J}$ and $\{B_k\}_{k\in K}$ be finite pairwise
    non-repeated families. Suppose every tensor in either family has a tensor
    in the other family whose mixed overlap does not tend to zero, and every
    such non-decaying pair is repeated. Then there is a bijection
    $\pi\colon J\to K$ such that $A_j$ and $B_{\pi(j)}$ are repeated for every
    $j\in J$.

    \emph{Scope restriction.} This result assumes the overlap consequences
    that Theorem~3.4 derives from proportionality of the assembled
    matrix-product vectors. It is the finite matching step, not the source
    theorem itself.
\end{theorem}
\begin{proof}\leanok
    Choose a non-decaying partner $f(j)\in K$ for every $j\in J$. If
    $f(j_1)=f(j_2)$, the two repeated-block relations compose to show that
    $A_{j_1}$ and $A_{j_2}$ are repeated, so pairwise non-repetition makes $f$
    injective. The symmetric construction gives an injection $K\to J$.
    Therefore $|J|=|K|$, and the first injection is the required bijection.
\end{proof}

\begin{theorem}[Proportional Fundamental Theorem for irreducible-form MPS]
    \label{thm:periodic_ft_proportional}
    \lean{MPSTensor.fundamentalTheorem_periodic_proportional_sectorDecomposition}
    \leanok
    \uses{def:nonzero_proportional_mpv, def:spectral_periodic_tensor}
    Let $A$ and $B$ have multiplicity-bearing irreducible forms
    \begin{align}
        A^i & =\bigoplus_{j\in J}(R_j\otimes A_j^i),
        \notag                                       \\
        B^i & =\bigoplus_{k\in K}(S_k\otimes B_k^i),
        \notag
    \end{align}
    where $R_j$ and $S_k$ are diagonal with nonzero diagonal entries and the
    two displayed block families are pairwise non-repeated bases of periodic
    tensors, spectrally periodic of periods $m_j$ and $n_k$ respectively.
    Suppose that, for every positive $N$, the vectors
    $\ket{V^{(N)}(A)}$ and $\ket{V^{(N)}(B)}$ are proportional by a nonzero
    scalar. Then there is a bijection $\pi\colon J\to K$ such that, for every
    $j\in J$, one has $m_j=n_{\pi(j)}$ and there
    are a unit-modulus scalar $\xi_j$ and an invertible matrix $Y_j$ such that,
    for every physical index $i$,
    \begin{align}
        A_j^i & =\xi_jY_jB_{\pi(j)}^iY_j^{-1}.
        \notag
    \end{align}
    Each $B_k$ occurs for exactly one $j$. This is
    \cite[Theorem~3.4, lines~613--623]{DeLasCuevas2017Irreducible}.

    The theorem relates the two bases of periodic tensors but makes no
    assertion relating their multiplicity matrices.
\end{theorem}
\begin{proof}
    \uses{thm:repeated_blocks_period_eq, thm:periodic_conditional_block_matching, thm:periodic_spectral_multiplicity_nondecaying_overlap}
    \leanok
    Theorem~\ref{thm:periodic_spectral_multiplicity_nondecaying_overlap}
    supplies non-decaying partners in both directions and classifies every
    such pair as repeated. Applying
    Theorem~\ref{thm:periodic_conditional_block_matching} gives a bijection
    $\pi\colon J\to K$ of repeated pairs. Finally,
    Theorem~\ref{thm:repeated_blocks_period_eq} gives $m_j=n_{\pi(j)}$ for every
    matched pair.
\end{proof}
